% compile 2 times
\documentclass[a4paper, 10pt]{extarticle}
\usepackage[italian]{babel}
\usepackage[
	solutions,%show/hide the solutions
	blue%color of solutions font
		    ]{verifiche}
\usepackage[margin=3cm]{geometry}
\usepackage{lipsum}
\usepackage{siunitx}


\begin{document}
%\renewcommand{\institutefont}{\huge}

\institute{Scuola superiore di Paperopoli}
\asyear[Anno accademico]{2019/2020}
\duration[Tempo:]{1 ora}
\testtype{Compito di fisica}
\instruction{Risolvere il seguente esercizio nel più breve tempo possibile.\\}
\printheading

\pagestyle{verifiche}

Nel presente documento è possibile vedere i risultati grafici delle varie tipologie di quesiti e gli altri strumenti forniti dal pacchetto \textsf{verifiche}. 

\begin{esercizio}
Risolvere il seguente problema.

Un mattone pesa \SI{1}{kg} più mezzo mattone. Quanto pesa un mattone?
\end{esercizio}

\begin{soluzione}
Il mattone pesa \SI{2}{kg}.
\end{soluzione}

\begin{esercizio}[pt=2]\label{es:ferropiuma}%
Pesa di più un chilogrammo di ferro o un chilogrammo di piume?
\end{esercizio}

\begin{esercizio}[diff=1]
Pesa di più un chilogrammo di ferro sulla Terra o un chilogrammo di piume su Giove?
\end{esercizio}

\renewcommand{\diffsymb}{$\bigstar$}
\begin{esercizio}[diff=2, pt=5]\label{es:grav}%
Scrivere la legge di gravitazione universale e commentarla.
\end{esercizio}

\begin{esercizio}[label=Esercizio]
Calcolare:
\[
\int_{-\infty}^{+\infty} e^{-x^2}\, dx
\]
\end{esercizio}

\begin{esercizio}
Rispondere alle seguenti domande:
\begin{enumerate}[(a)]
\item \partialpt{2} Descrivere la legge di Boyle
\item \partialpt{3} Rappresentarla graficamente
\end{enumerate}
\end{esercizio}

\begin{esercizio}[diff=3, partialpt]\label{es:meccanica}
Rispondere alle seguenti domande:
\begin{enumerate}[(i)]
\item\partialpt{2} Descrivere l'energia meccanica di un sistema gravitazionale a due corpi isolato.
\item\partialpt{3} Enunciare il teorema del viriale.
\end{enumerate}
\end{esercizio}

\begin{esercizio*}[label=Domanda, pt=2, diff=1]
Calcolare
\[
\frac{d}{dx}\left[ x^2\right]
\]
\end{esercizio*}

\setlength{\ptrulerlength}{.5cm}
\begin{esercizio}[pt=2]
Calcolare
\[
\frac{d}{dx}\left[ x^2\right]
\]
\end{esercizio}

\ptprefix{.../}
\ptdelimiters{[]}
\ptlabel{punti}
\begin{esercizio}[pt=2]\label{es:derivata}
Calcolare la seguente derivata:
\[
\frac{d e^x}{dx} = 
\]
\end{esercizio}

\partialptprefix{.../}
\partialptdelimiters{{}{}}
\partialptlabel[punto]{punti}
\begin{esercizio}
Risolvere i seguenti quesiti:
\begin{enumerate}[(i)]
\item\partialpt{2}\noindent Calcolare il seguente integrale indefinito:
\[
\int e^x dx = 
\]
\item\partialpt{1}\noindent Calcolare il seguente integrale definito:
\[
\int_0^1 e^x dx = 
\]
\end{enumerate}
\end{esercizio}


\begin{soluzione}
Questa è la soluzione dell'esercizio \ref{es:ferropiuma} a pagina \pageref{es:ferropiuma}.
\end{soluzione}

\begin{soluzione}[label={Soluzione dell'esercizio \ref{es:ferropiuma}}]
Soluzione dell'esercizio a pagina \pageref{es:ferropiuma}.
\end{soluzione}



\begin{esercizio}
La branca della fisica che si occupa dello studio di \emph{come} si muovono i corpi è detta: \inlinesol{cinematica.}
\end{esercizio}

\begin{esercizio}
La \completetext{Terra} è il pianeta su cui viviamo e la sua unica \completetext{stella} è il Sole.
Il pianeta \completetext{Giove} è il secondo corpo celeste, per dimensione, dopo il Sole.
\end{esercizio}

\begin{esercizio}
Indicare se le seguenti affermazioni sono vere o false.\par
\begin{tabular}{p{.7\textwidth}l}
Tutti gli ateniesi mentono, ed io sono ateniese	&	\truefalse{}\\[.5em]
Il Sole è un pianeta del Sistema solare		&	\truefalse{F}\\[.5em]
Io sono vera							&	\truefalse{V}
\end{tabular}
\end{esercizio}

\begin{esercizio}
Quali delle seguenti equazioni descrive il moto rettilineo uniforme?
\begin{closedquestion}
\item $a=\SI{5}{\m\per\s\squared}$
\item[\checked] $x(t) = v t + x_0$
\item $x(t) = t^3$
\item $t = a^2$
\end{closedquestion}
\end{esercizio}

\renewcommand{\closedquestionitem}{$\bigcirc$}
\renewcommand{\checkmarker}{$\times$\,}
\begin{esercizio}
Quali delle seguenti equazioni descrive il moto rettilineo uniforme?\par
\begin{closedquestion*}
\item $a=\SI{5}{\m\per\s\squared}$
\item[\checked] $x(t) = v t + x_0$
\item $x(t) = t^3$
\item $t = a^2$
\end{closedquestion*}
\end{esercizio}

\begin{esercizio}
\lipsum[2]\\[1em]
\openquestion[height=5cm]{\lipsum[1]}
\end{esercizio}

\begin{esercizio}
\lipsum[2]\\[1em]
\openquestion[type=lines, width=8cm, height=6cm, linecolor=orange]{\lipsum[2]}
\end{esercizio}


\begin{esercizio}
Trovare gli errori nel seguente testo e correggerli.

\renewcommand{\baselinestretch}{2.0}
\sffamily
I numeri preceduti dal segno "$+$" o dal segno "$-$" si dicono \finderror{numeri razionali}{numeri relativi} in ragione del segno che li precede.
Nel prodotto di due numeri con segno, la moltiplicazione di due numeri opposti è un numero \finderror{positivo}{negativo}.
A differenza dei numeri assoluti, la differenza tra due  \finderror{numeri razionali}{numeri relativi} con segno \finderror{non è sempre possibile}{è sempre possibile} (ad esempio $5 - 7$). 
\renewcommand{\baselinestretch}{1.0}
\end{esercizio}
\baselinestretch

\begin{esercizio}
\textandimage{Dato il seguente grafico indicare sul grafico i punti di massimo globale.}{
\begin{tikzpicture}
\draw[latex-latex] (3,0) node[right] {$x$} -- (0,0) -- (0,3) node[above]{$y$};
\draw plot[domain=0:3, samples=100] (\x, {2.5*exp(-\x^2)});
\inlinesol{\draw[fill, red] (0,2.5) circle (2pt);}
\end{tikzpicture}
}
\end{esercizio}

\begin{esercizio}[pt=3]\label{es:ultimo}
\makecolumn{.5\textwidth}{%
	Dato il seguente grafico indicare sul grafico i punti di massimo globale.
}
\makecolumn{.4\textwidth}{
\hspace*{\fill}
\fbox{
	\begin{tikzpicture}[baseline=(current bounding box.north)]
	\draw[latex-latex] (3,0) node[right] {$x$} -- (0,0) -- (0,3) node[above]{$y$};
	\draw plot[domain=0:3, samples=100] (\x, {2.5*exp(-\x^2)});
	\inlinesol{\draw[fill, red] (0,2.5) circle (2pt);}
	\end{tikzpicture}}
\hspace*{\fill}
\\[1em]
\lipsum[2]
}
\end{esercizio}


\vfill
\begin{center}
\begin{tabular}{|c|c|c|c|c|c|c|}
\hline
Esercizio	&	Quesito \ref{es:ferropiuma}
		&	Quesito \ref{es:grav}
		&	Quesito \ref{es:meccanica}
		&	Quesito \ref{es:derivata}
		&	Quesito \ref{es:ultimo}
		&	Totale
		\\
 		&	\ref{ptes@2} pt
		&	\ref{ptes@4} pt
		&	\ref{ptes@7} pt
		&	\ref{ptes@9} pt
		&	\ref{ptes@20} pt
		&	
		\\
\hline
Punteggio	&&&&&&\\[1em]
\hline
\multicolumn{4}{c}{}&&Voto&\\[1em]
\cline{6-7}
\end{tabular}
\end{center}


\clearpage
\pagestyle{empty}
\vspace*{\fill}
\begin{center}
\huge Esempio \verb|multitest|
\end{center}
\vspace*{\fill}
\clearpage

%\pagestyle{plain}
\pagestyle{mainverifiche}
\institute{Scuola superiore di Paperopoli}
\asyear[Anno accademico]{2023/2024}
\duration[Tempo:]{2 ora}
\testtype{Compito di matematica}
\instruction{Risolvere il seguente esercizio nel più breve tempo possibile.\\}

\pgfkeys{/pgf/number format/.cd,
	%zerofill=true,
	std,
	%int detect,
	precision=3,
	%set decimal separator={$,$},
	use comma,
	set thousands separator={$\,$}
}


\renewcommand{\ptfont}{\footnotesize}

\begin{multitest}[4]
%\pgfmathsetseed{\themultitestcounter*42}
\subtitle{Test \Alph{multitestcounter}}
\headingstyle{einstein}
\printheading
\pgfmathrandominteger{\ax}{-2}{2}
\pgfmathrandominteger{\ay}{-2}{2}
\pgfmathrandominteger{\bx}{-2}{2}
\pgfmathrandominteger{\by}{-2}{2}

\begin{esercizio}[partialpt]
\partialpt{3}	Disegna, in un grafico cartesiano, i seguenti punti:
	\[
	A = (\ax; \ay) \quad B = (\bx; \by)
	\]
\partialpt{5}	e calcola la loro distanza.
\end{esercizio}

\begin{soluzione}
	Si rappresentano i punti dati nel piano cartesiano
	\begin{center}
	\begin{tikzpicture}
	\draw[gray] (-3,-3) grid (3,3);
	\draw[->, thick] (-3,0) -- (3,0) node[right]{x};
	\draw[->, thick] (0,-3) -- (0,3) node[above]{y};
	\draw[red, very thick] (\ax, \ay) -- (\bx, \by);
	\fill (\ax, \ay) circle (2pt) node[above] {A};
	\fill (\bx, \by) circle (2pt) node[below] {B};
	\end{tikzpicture}
	\end{center}
La distanza tra i due punti è data dall'equazione:
\[
d(A, B)= \sqrt{(x_A - x_B) ^2 + (y_A- y_B)^2} = \pgfmathparse{sqrt((\ax-\bx)^2+(\ay-\by)^2)}\pgfmathprintnumber\pgfmathresult
\]
\end{soluzione}
\end{multitest}

\begin{esercizio}
	Per ogni domanda indicare quale delle risposte è quella corretta.
\begin{shuffledenumerate}[(i)]
\sitem{\item Il corpo rigido è:
	\begin{shuffledclosed}
		\sitem{\item un oggetto assolutamente non elastico.}
		\sitem{\item un oggetto esteso che subisce deformazioni quando gli vengono applicate forze molto intense.}
		\sitem{\item[\checked] un oggetto esteso che non subisce deformazioni qualsiasi siano le forze che gli vengono applicate.}
		\sitem{\item un modello che descrive il comportamento del punto materiale.}
	\end{shuffledclosed}
}
\sitem{\item È possibile mantenere in equilibrio un ombrello sulla punta di un dito solo se:
	\begin{shuffledclosed}
	\sitem{\item il peso dell’ombrello è equilibrato dalla forza esercitata dal dito.}
	\sitem{\item[\checked] il baricentro dell’ombrello sta sulla retta verticale che passa per la punta del dito.}
	\sitem{\item il momento totale agente sull’ombrello è positivo.}
	\sitem{\item la forza totale agente sull'ombrello è positivo.}
	\end{shuffledclosed}
}

\sitem{\item La condizione di equilibrio per un punto materiale è che:
	\begin{shuffledclosed}
	\sitem{\item non ci sia alcuna forza agente su di esso.}
	\sitem{\item ci siano soltanto forze vincolari agenti su di esso.}
	\sitem{\item[\checked] la somma delle forze ad esso applicate sia uguale a zero.}
	\sitem{\item la somma delle forze vincolari agenti su di esso sia uguale a zero.}
	\end{shuffledclosed}
}

\sitem{\item Quale delle seguenti definizioni è corretta?
	\begin{shuffledclosed}
	\sitem{\item  Il baricentro di un oggetto qualunque è il centro geometrico dell'oggetto.}
	\sitem{\item Il centro di simmetria di un oggetto qualunque è il centro geometrico dell'oggetto.}
	\sitem{\item[\checked] II baricentro di un oggetto qualunque è il punto di applicazione della forza-peso totale agente sull'oggetto.}
	\sitem{\item Il centro di simmetria di un oggetto qualunque è il punto di applicazione della forza-peso totale agente sull'oggetto.}
	\end{shuffledclosed}
}
\end{shuffledenumerate}
\end{esercizio}


\end{document}